% CP-VI-0018
% target_chapter: VI/16 — Exact Sequences of Sheaves
% source_unit_id: IMP-DL-0E3E7EC766-P018
% source: Downloads/theory-of-algebraic-geometry-1.tex L3821-L3849
% solution_provenance: SOURCE_EXPLICIT_SOLUTION

\subsection*{Exercise 6}

Show that a sequence of sheaves
\[
\cdots
\longrightarrow
\mathcal F_{i-1}
\xrightarrow{d_{i-1}}
\mathcal F_i
\xrightarrow{d_i}
\mathcal F_{i+1}
\longrightarrow
\cdots
\]
is exact if and only if for every $p \in X$, the corresponding sequence of
stalks
\[
\cdots
\longrightarrow
(\mathcal F_{i-1})_p
\xrightarrow{(d_{i-1})_p}
(\mathcal F_i)_p
\xrightarrow{(d_i)_p}
(\mathcal F_{i+1})_p
\longrightarrow
\cdots
\]
is exact.

\subsection*{Solution}

We prove that exactness of sheaves is a local property.

Recall that the sequence is exact at $\mathcal F_i$ if
\[
\operatorname{im}(d_{i-1})
=
\ker(d_i)
\]
as subsheaves of $\mathcal F_i$.

We must show that this happens if and only if
\[
\operatorname{im}\bigl((d_{i-1})_p\bigr)
=
\ker\bigl((d_i)_p\bigr)
\]
for every point $p \in X$.

\subsection*{1. Exactness of sheaves implies exactness of stalks}

Assume that the sequence of sheaves is exact.

Thus, for every $i$,
\[
\operatorname{im}(d_{i-1})
=
\ker(d_i)
\]
as sheaves.

Taking stalks at $p$, we get
\[
\bigl(\operatorname{im}(d_{i-1})\bigr)_p
=
\bigl(\ker(d_i)\bigr)_p.
\]

Now stalks commute with kernels:
\[
\bigl(\ker(d_i)\bigr)_p
=
\ker\bigl((d_i)_p\bigr).
\]

Also, by the description of sheaf images from Exercise 5,
\[
\bigl(\operatorname{im}(d_{i-1})\bigr)_p
=
\operatorname{im}\bigl((d_{i-1})_p\bigr).
\]

Therefore
\[
\operatorname{im}\bigl((d_{i-1})_p\bigr)
=
\ker\bigl((d_i)_p\bigr).
\]

Hence the stalk sequence is exact at $(\mathcal F_i)_p$.

Since $i$ and $p$ were arbitrary, the sequence of stalks is exact at every
point.

\subsection*{2. Exactness of stalks implies exactness of sheaves}

Now assume that for every $p \in X$, the sequence of stalks is exact.

Thus, for every $p$,
\[
\operatorname{im}\bigl((d_{i-1})_p\bigr)
=
\ker\bigl((d_i)_p\bigr).
\]

We must prove
\[
\operatorname{im}(d_{i-1})
=
\ker(d_i)
\]
as sheaves.

First observe that
\[
d_i \circ d_{i-1} = 0
\]
as a morphism of sheaves.

Indeed, on stalks,
\[
(d_i)_p \circ (d_{i-1})_p = 0
\]
for every $p$ by exactness of the stalk sequence. Since a morphism of sheaves
is zero if and only if it is zero on every stalk, it follows that
\[
d_i \circ d_{i-1} = 0.
\]

Therefore
\[
\operatorname{im}(d_{i-1})
\subseteq
\ker(d_i).
\]

We need to prove the reverse inclusion.

Let
\[
s \in \ker(d_i)(U)
\]
for some open set $U \subseteq X$.

This means
\[
s \in \mathcal F_i(U)
\]
and
\[
d_i(s)=0
\quad \text{in } \mathcal F_{i+1}(U).
\]

We must show that $s$ is locally in the image of $d_{i-1}$.

Let $p \in U$.

The germ
\[
s_p \in (\mathcal F_i)_p
\]
satisfies
\[
(d_i)_p(s_p)=0.
\]
Thus
\[
s_p \in \ker\bigl((d_i)_p\bigr).
\]

By exactness of the stalk sequence,
\[
s_p \in \operatorname{im}\bigl((d_{i-1})_p\bigr).
\]

Therefore there exists a germ
\[
t_p \in (\mathcal F_{i-1})_p
\]
such that
\[
(d_{i-1})_p(t_p)=s_p.
\]

Choose an open neighborhood
\[
V \subseteq U
\]
of $p$ and a section
\[
t \in \mathcal F_{i-1}(V)
\]
representing the germ $t_p$.

Then
\[
(d_{i-1}(t))_p = s_p.
\]

Equality of germs means that after possibly shrinking $V$, we have
\[
d_{i-1}(t) = s|_V.
\]

Thus, for every point $p \in U$, there exists a neighborhood $V$ of $p$ such
that
\[
s|_V \in \operatorname{im}(d_{i-1,V}).
\]

Therefore $s$ is locally in the image of $d_{i-1}$.

By the description of the sheaf image,
\[
s \in \operatorname{im}(d_{i-1})(U).
\]

Hence
\[
\ker(d_i)(U)
\subseteq
\operatorname{im}(d_{i-1})(U)
\]
for every open set $U$.

Thus
\[
\ker(d_i)
\subseteq
\operatorname{im}(d_{i-1}).
\]

Together with the opposite inclusion, we obtain
\[
\ker(d_i)
=
\operatorname{im}(d_{i-1}).
\]

Therefore the sequence of sheaves is exact.

\subsection*{Conclusion}

We have shown both implications:

\[
\boxed{
	\text{A sequence of sheaves is exact}
	\iff
	\text{the induced sequence of stalks is exact at every point.}
}
\]

Equivalently, exactness of sheaves is a pointwise condition.

The reason is that sheaves are local objects: a section belongs to a kernel or
to an image sheaf precisely when the corresponding statement holds locally,
and local statements are detected by germs.
